% sections/07_overview_viz.tex — V1, V2, V3, V4

\section{Overview Visualisations}

%% ── V1 ──────────────────────────────────────────────────────
\subsection{V1: Geographic Distribution of High-Severity Accidents}

The dot map plots each accident by GPS coordinates, with Severity~3 in blue and
Severity~4 in red. Low-severity accidents (1--2) are excluded to highlight the
spatial concentration of disruptive incidents.

% TO ADD FIGURE: save chart as figures/v1_geo_map.png then replace the
% placeholder block below with:
 % \figchart{v1_geo_map}{Geographic distribution of high-severity accidents}{fig:v1_geo_map.png}
\begin{figure}[H]
  \centering
   \figchart{v1_geo_map}{Geographic distribution of high-severity accidents (Severity 3--4)}{fig:v1_geo_map.png}
  \label{fig:v1}
\end{figure}

\idiomtable
  {0D Point}
  {Position (ordered, unaligned) + Colour Hue (blue = Sev.\,3, red = Sev.\,4)}
  {Identify spatial pattern in accident severity distribution}
  {500K points; overplotting managed by point size and opacity}

\insightbox{Severity 4 accidents (red) cluster densely along the Eastern Seaboard
and major interstate corridors, particularly in Florida, California, and the Texas
triangle. The western interior shows predominantly Severity~3 (blue) in sparse
highway networks. The geographic skew mirrors population density and road network
complexity rather than environmental conditions alone.}

%% ── V2 ──────────────────────────────────────────────────────
\subsection{V2: Top 25 States by Total Accident Count (2016--2022)}

A ranked horizontal bar chart showing California leading with approximately
1.39M accidents, more than double Florida (687K) and three times Texas (454K).
The sorted order enables immediate rank lookup and magnitude comparison.

% TO ADD FIGURE: save as figures/v2_top25_states.png
% \figchart{v2_top25_states}{Top 25 states by total accident count}{fig:v2}
\begin{figure}[H]
  \centering
  \figchart{v2_top25_states}{Top 25 states by total accident count}{fig:v2}
  \label{fig:v2}
\end{figure}

\idiomtable
  {1D Line (bar)}
  {Horizontal Length (ordered, aligned)}
  {Lookup rank and compare magnitude across states}
  {Dozens of categories, thousands of values per bar}

\insightbox{California's dominance is partly a data-collection artefact --- the
dataset relies on traffic API feeds, and California has one of the most
comprehensive sensor networks in the US.}

%% ── V3 ──────────────────────────────────────────────────────
\subsection{V3: Overall Severity Distribution}

A simple bar chart reveals the sharp skew: Severity~2 dominates with 4.85M
records, while Severity~1 (65K), Severity~3 (1.15M), and Severity~4 (103K) are
comparatively small. This distribution is critical for interpreting all
downstream High Severity Rate calculations.

% TO ADD FIGURE: save as figures/v3_severity_dist.png
% \figchart{v3_severity_dist}{Overall severity distribution}{fig:v3}
\begin{figure}[H]
  \centering
  \figchart{v3_severity_dist}{Overall severity distribution across the full dataset}{fig:v3}
  \label{fig:v3}
\end{figure}

\idiomtable
  {1D Bar}
  {Vertical Position + Colour Hue (categorical identity)}
  {Compare counts across ordinal severity categories}
  {4 categories --- optimal for grouped bars}

\insightbox{Over 80\% of accidents are Severity~2. This class imbalance means
High Severity Ratio (proportion of Severity 3--4) is a more meaningful signal
than raw counts throughout this report. A shift from 19\% to 22\% high severity
represents a substantial change in absolute terms given the volume.}

%% ── V4 ──────────────────────────────────────────────────────
\subsection{V4: Year-wise Accident Volume and High Severity Rate}

A dual-axis combo chart pairs stacked bars (High vs.\ Low severity counts) with
a dashed line for the average yearly severity rate. The COVID-19 period
(2020--2021) is highlighted in pink.

% TO ADD FIGURE: save as figures/v4_yearly_trend.png
% \figchart{v4_yearly_trend}{Year-wise volume and high-severity rate}{fig:v4}
\begin{figure}[H]
  \centering
  \figchart{v4_yearly_trend}{Year-wise accident volume (stacked bars) and high-severity rate (dashed line)}{fig:v4}
  \label{fig:v4}
\end{figure}

\idiomtable
  {1D Bar (stacked) + 0D Point}
  {Vertical Length + Colour Hue (stacked) + Position on secondary Y-axis (rate)}
  {Compare trend over time; identify COVID anomaly}
  {7 years (2016--2022) --- manageable temporal resolution}

\insightbox{High severity rate declined sharply from 34\% in 2016 to under 3\%
by 2022. This counter-intuitive long-term trend likely reflects the rapid growth
in Severity~2 records as API coverage expanded year-on-year, diluting the
high-severity proportion --- not a genuine improvement in accident outcomes.
The COVID dip in 2020 shows reduced accident volume but the rate continued its
structural decline. Inter-year severity rate comparisons should be treated as
exploratory, not conclusive.}
